\documentclass[11pt]{article}

% ---- Packages ----
\usepackage{amsmath, amssymb, amsthm}
\usepackage{mathtools}
\usepackage{bm}
\usepackage[margin=1in]{geometry}
\usepackage{hyperref}
\usepackage{booktabs}
\usepackage{array}
\usepackage{enumitem}
\usepackage{listings}
\usepackage{xcolor}

\hypersetup{
  colorlinks=true,
  linkcolor=blue!50!black,
  citecolor=blue!50!black,
  urlcolor=blue!50!black
}

% ---- Theorem environments ----
\newtheorem{proposition}{Proposition}
\newtheorem{lemma}{Lemma}
\theoremstyle{definition}
\newtheorem{definition}{Definition}
\newtheorem{remark}{Remark}

% ---- Notation shortcuts ----
\newcommand{\R}{\mathbb{R}}
\newcommand{\Rplus}{\mathbb{R}_{+}}
\newcommand{\E}{\mathbb{E}}
\newcommand{\Var}{\operatorname{Var}}
\newcommand{\Cov}{\operatorname{Cov}}
\newcommand{\Norm}{\mathcal{N}}
\newcommand{\Unif}{\mathrm{Unif}}
\newcommand{\Gam}{\mathrm{Gamma}}
\newcommand{\HN}{\mathrm{HN}}
\newcommand{\given}{\,\vert\,}
\newcommand{\Phinv}{\Phi^{-1}}
\newcommand{\code}[1]{\texttt{#1}}

% ---- Listings setup for R code ----
\lstset{
  basicstyle=\ttfamily\footnotesize,
  breaklines=true,
  showstringspaces=false,
  language=R,
  commentstyle=\color{gray},
  keywordstyle=\color{blue!60!black},
  stringstyle=\color{green!50!black},
  frame=single,
  framerule=0.3pt
}

% ---- Title block ----
\title{
  Copula / Probit Transforms and the Autocorrelated Gaussian Prior\\
  for Time-Varying Parameters in the Spatial Skew-$t$ Model:\\
  Code-to-Text Correspondence
}
\author{Technical Note --- \texttt{transform\_ar\_construction}}
\date{}

\begin{document}
\maketitle

\begin{abstract}
We document and verify the construction by which the time-varying parameters of
the spatial skew-$t$ process acquire their autocorrelation. The governing idea is
stated in one sentence:
\begin{quote}
\emph{``For each time-varying parameter, we transform the parameter to obtain a
standard normal marginal distribution, place a Gaussian prior with autocorrelation
on the transformed parameter, and then transform back to the appropriate marginal
distribution for the skew-$t$ process.''}
\end{quote}
We locate the three operations of this recipe in the \code{ar2} source tree, give
the exact mathematical form of each transform and its inverse CDF, and confirm
that the implementation is faithful to the text. The key implementation detail
that makes the statement true---namely that the autocorrelated Gaussian prior is
normalised to unit marginal variance at \emph{every} time point so the
standard-normal marginal is preserved---is made explicit.
\end{abstract}

% ====================================================================
\section{The three-step recipe and where it lives in the code}
\label{sec:recipe}

There are three time-varying parameters, each indexed by knot $k$ and day
$t\in\{1,\dots,n_t\}$:
\begin{itemize}[nosep]
  \item the partition precision $\tau_{kt}$ (\code{tau});
  \item the skew random effect $z_{kt}$ (\code{z});
  \item the knot location $w_{kt}=(w_{kt}^{(1)},w_{kt}^{(2)})$ (\code{knots}).
\end{itemize}
Each parameter follows the identical three-step recipe. Table~\ref{tab:map}
maps the sentence onto the source files.

\begin{table}[h]
\centering
\renewcommand{\arraystretch}{1.3}
\begin{tabular}{p{0.30\textwidth} p{0.38\textwidth} p{0.24\textwidth}}
\toprule
\textbf{Text concept} & \textbf{Code operation} & \textbf{Location} \\
\midrule
\textcircled{1} transform $\to$ standard normal marginal
  & forward copula / probit: \code{gamma.cop}, \code{hn.cop}, \code{transform\$probit}
  & \code{auxfunctions.R} L25--33, L78--81 \\
\textcircled{2} Gaussian prior with autocorrelation
  & stationary AR(2) conditional normal densities, \code{ar2\_conditional\_log\_density}
  & \code{auxfunctions\_ar2.R} L56--102 \\
\textcircled{3} transform back $\to$ skew-$t$ marginal
  & inverse copula / inv.\ probit: \code{gamma.invcop}, \code{hn.invcop}, \code{transform\$inv.probit}
  & \code{auxfunctions.R} L29--33 \\
\bottomrule
\end{tabular}
\caption{Correspondence between the textual recipe and the implementation.
All paths are relative to \code{code/R/ar2/}.}
\label{tab:map}
\end{table}

The Metropolis updates that drive these parameters live in
\code{update\_params\_ar2.R}: \code{updateTauTS\_AR2}, \code{updateZTS\_AR2}, and
\code{updateKnotsTS\_AR2}. In each, the proposal and acceptance are carried out
\emph{on the transformed (Gaussian) scale}, which is natural because the prior is
Gaussian there.

% ====================================================================
\section{The transform tools}
\label{sec:tools}

Throughout, $\Phi$ and $\varphi$ denote the standard normal CDF and PDF, and
$\Phinv=\code{qnorm}$.

\begin{definition}[Gaussian copula transform]
For a target marginal CDF $F$, the object \code{transform\$copula}
(\code{auxfunctions.R} L25--28) and its inverse \code{transform\$inv.copula}
(L29--32) implement
\begin{equation}
\operatorname{cop}_F(x) = \Phinv\!\big(F(x)\big),
\qquad
\operatorname{cop}_F^{-1}(u^\ast) = F^{-1}\!\big(\Phi(u^\ast)\big).
\label{eq:copula}
\end{equation}
\end{definition}

\begin{definition}[Probit transform]
For $\Unif(l,u)$, \code{transform\$probit} (L14--17) and \code{transform\$inv.probit}
(L18--22) implement
\begin{equation}
\operatorname{probit}(w) = \Phinv\!\Big(\tfrac{w-l}{u-l}\Big),
\qquad
\operatorname{probit}^{-1}(w^\ast) = (u-l)\,\Phi(w^\ast) + l.
\label{eq:probit}
\end{equation}
\end{definition}

The probit pair is the copula special case of \eqref{eq:copula} with
$F(w)=(w-l)/(u-l)$, the $\Unif(l,u)$ CDF. The named copulas instantiated in
\code{auxfunctions.R} (L78--81) are
\[
\code{gamma.cop} = \operatorname{cop}_{F_\Gam},\quad
\code{gamma.invcop} = \operatorname{cop}_{F_\Gam}^{-1},\quad
\code{hn.cop} = \operatorname{cop}_{F_\HN},\quad
\code{hn.invcop} = \operatorname{cop}_{F_\HN}^{-1}.
\]

\begin{lemma}[Standard-normal marginal]
\label{lem:sn}
If $X$ has continuous marginal CDF $F$, then $\operatorname{cop}_F(X)=\Phinv(F(X))\sim\Norm(0,1)$,
since $F(X)\sim\Unif(0,1)$ and $\Phinv$ maps $\Unif(0,1)$ to $\Norm(0,1)$.
Conversely $\operatorname{cop}_F^{-1}(U^\ast)=F^{-1}(\Phi(U^\ast))$ has marginal $F$
whenever $U^\ast\sim\Norm(0,1)$.
\end{lemma}

Lemma~\ref{lem:sn} is exactly step~\textcircled{1} (forward) and step~\textcircled{3}
(backward) of the recipe. It also pins down the requirement that makes the recipe
correct: in step~\textcircled{3} the back-transform reproduces the intended marginal
$F$ \emph{only if} the transformed variate is $\Norm(0,1)$. Section~\ref{sec:prior}
shows the AR(2) prior is engineered to guarantee this at every $t$.

% ====================================================================
\section{Per-parameter marginals, forward transforms, and inverse CDFs}
\label{sec:perparam}

\subsection{Precision $\tau$: Gamma marginal}
The marginal is $\tau\sim\Gam(\alpha,\beta)$ with shape $\alpha=\code{tau.alpha}$
and rate $\beta=\code{tau.beta}$.

\paragraph{Forward (to $\Norm(0,1)$).}
\code{update\_params\_ar2.R} L19, \code{tau.star <- gamma.cop(tau, tau.alpha, tau.beta)}:
\begin{equation}
\tau^\ast = \Phinv\!\big(F_\Gam(\tau;\alpha,\beta)\big),
\qquad
F_\Gam(\tau;\alpha,\beta) = \frac{\gamma(\alpha,\beta\tau)}{\Gamma(\alpha)},
\end{equation}
where $\gamma(\cdot,\cdot)$ is the lower incomplete gamma function.

\paragraph{Inverse CDF (back to $\Rplus$).}
\code{update\_params\_ar2.R} L31, \code{can.tau <- gamma.invcop(can.tau.star, tau.alpha, tau.beta)}:
\begin{equation}
\tau = F_\Gam^{-1}\!\big(\Phi(\tau^\ast);\alpha,\beta\big)
     = Q_\Gam\!\big(\Phi(\tau^\ast)\big),
\end{equation}
with $Q_\Gam=\code{qgamma}$ the gamma quantile function. A floor of $10^{-6}$ is
applied for numerical stability (L32--34).

\subsection{Skew effect $z$: Half-Normal marginal}
Conditional on $\tau$, the model takes $z\given\tau\sim\HN(0,\sigma)$ with scale
$\sigma=1/\sqrt{\tau}$ (\code{update\_params\_ar2.R} L343--344). The half-normal CDF and
quantile function are (\code{auxfunctions.R} L56--73)
\begin{equation}
F_\HN(z;\sigma) = 2\Phi\!\big(z/\sigma\big) - 1,
\qquad
F_\HN^{-1}(p;\sigma) = \sigma\,\Phinv\!\Big(\tfrac{1+p}{2}\Big),
\qquad z \ge 0.
\end{equation}

\paragraph{Forward (to $\Norm(0,1)$).}
\code{z.star <- hn.cop(z, sig)}:
\begin{equation}
z^\ast = \Phinv\!\big(2\Phi(z\sqrt{\tau}) - 1\big).
\end{equation}

\paragraph{Inverse CDF (back to $\Rplus$).}
\code{update\_params\_ar2.R} L372, \code{can.z <- hn.invcop(can.z.star, sig)}:
\begin{equation}
z = \frac{1}{\sqrt{\tau}}\;\Phinv\!\Big(\tfrac{1+\Phi(z^\ast)}{2}\Big).
\end{equation}

\begin{remark}[A latent inconsistency in \code{dhn}, off the transform path]
The density helper \code{dhn} returns \code{dnorm(x,0,sig)/2} (\code{auxfunctions.R} L47),
i.e.\ $\varphi(x/\sigma)/(2\sigma)$, whereas the true half-normal density consistent
with $F_\HN$ above is $2\varphi(x/\sigma)/\sigma$---a factor of four apart. The copula
transforms call only \code{phn} and \code{qhn}, so \code{dhn} never enters the AR
construction; the discrepancy is a latent bug with no effect on the pipeline.
\end{remark}

\subsection{Knot location $w$: Uniform marginal}
Each coordinate $j\in\{1,2\}$ has marginal $w^{(j)}\sim\Unif(l_j,u_j)$ with
$l_j=\code{min.s[j]}$, $u_j=\code{max.s[j]}$.

\paragraph{Forward (to $\Norm(0,1)$).}
\code{update\_params\_ar2.R} L477--484, \code{transform\$probit}:
\begin{equation}
w^{(j)\ast} = \Phinv\!\Big(\tfrac{w^{(j)} - l_j}{u_j - l_j}\Big).
\end{equation}

\paragraph{Inverse CDF (back to the spatial domain).}
\code{update\_params\_ar2.R} L520--527, \code{transform\$inv.probit}:
\begin{equation}
w^{(j)} = (u_j - l_j)\,\Phi(w^{(j)\ast}) + l_j.
\end{equation}

% ====================================================================
\section{Step \textcircled{2}: the autocorrelated Gaussian prior}
\label{sec:prior}

Write $X^\ast_t$ for the transformed series, i.e.\ any of $\tau^\ast_{kt}$,
$z^\ast_{kt}$, $w^{(j)\ast}_{kt}$. The prior is a stationary, mean-zero,
\emph{unit-variance} AR(2) (\code{auxfunctions\_ar2.R} L9--54):
\begin{equation}
X^\ast_t = \phi_1 X^\ast_{t-1} + \phi_2 X^\ast_{t-2} + \varepsilon_t,
\qquad \varepsilon_t \sim \Norm(0,\sigma_\varepsilon^2).
\label{eq:ar2}
\end{equation}
The decisive design choice is the normalisation $\gamma_0\equiv 1$
(\code{auxfunctions\_ar2.R} L31). The remaining stationary autocovariances and the
innovation variance then follow from the Yule--Walker equations:
\begin{equation}
\gamma_0 = 1,\qquad
\gamma_1 = \frac{\phi_1}{1-\phi_2},\qquad
\gamma_2 = \phi_1\gamma_1 + \phi_2,\qquad
\sigma_\varepsilon^2 = \gamma_0 - \phi_1\gamma_1 - \phi_2\gamma_2.
\label{eq:yw}
\end{equation}
Used as a prior, the joint density factorises into the conditional normals
(\code{auxfunctions\_ar2.R} L62--92):
\begin{align}
X^\ast_1 &\sim \Norm(0,\,1), \label{eq:cond1}\\
X^\ast_2 \given X^\ast_1 &\sim \Norm\!\big(\gamma_1 X^\ast_1,\; 1-\gamma_1^2\big), \label{eq:cond2}\\
X^\ast_t \given X^\ast_{t-1},X^\ast_{t-2} &\sim \Norm\!\big(\phi_1 X^\ast_{t-1} + \phi_2 X^\ast_{t-2},\; \sigma_\varepsilon^2\big), \quad t\ge 3. \label{eq:condt}
\end{align}
Stationarity is enforced by the triangle (\code{check\_ar2\_stability}, L2--7):
\begin{equation}
|\phi_2| < 1,\qquad \phi_1 + \phi_2 < 1,\qquad \phi_2 - \phi_1 < 1.
\label{eq:stab}
\end{equation}

\begin{proposition}[Marginal invariance]
\label{prop:marg}
Under \eqref{eq:ar2}--\eqref{eq:condt}, every $X^\ast_t$ has marginal $\Norm(0,1)$.
Consequently, by Lemma~\ref{lem:sn}, the back-transform of step~\textcircled{3}
returns each parameter to its intended marginal ($\Gam$, $\HN$, or $\Unif$) at every
time point.
\end{proposition}

\begin{proof}[Sketch]
Equation \eqref{eq:cond1} gives $\Var(X^\ast_1)=\gamma_0=1$. For \eqref{eq:cond2},
$\Var(X^\ast_2)=\gamma_1^2\Var(X^\ast_1)+(1-\gamma_1^2)=1$. For $t\ge 3$, the
Yule--Walker choice \eqref{eq:yw} of $\sigma_\varepsilon^2$ is precisely the value
that propagates $\Var(X^\ast_t)=1$ forward. Hence the marginal variance is $1$ for
all $t$; with zero mean this is $\Norm(0,1)$.
\end{proof}

% ====================================================================
\section{Verdict: the code is faithful to the text}
\label{sec:verdict}

The implementation realises each clause of the sentence, and additionally handles
two requirements that the prose leaves implicit but that correctness demands.

\begin{enumerate}[label=(\roman*), leftmargin=*]
\item \textbf{``standard normal marginal'' must hold at every $t$.}
The back-transform $F^{-1}(\Phi(X^\ast_t))$ recovers the correct $F$ only when
$X^\ast_t\sim\Norm(0,1)$. The $\gamma_0=1$ normalisation
(Proposition~\ref{prop:marg}) secures this. In particular, the $t=2$ conditional
\eqref{eq:cond2} uses the stationary bivariate form $\Norm(\gamma_1 X^\ast_1,1-\gamma_1^2)$
rather than the backward-extrapolated $\Norm(\phi_1 X^\ast_1,\sigma_\varepsilon^2)$; the
two coincide only when $\phi_2=0$, and only the former keeps $\Var=1$. This is the
``Fix~1'' branch in \code{get\_ar2\_conditional\_params} (\code{auxfunctions\_ar2.R}
L70--85).

\item \textbf{The Metropolis ratio must include the $t{+}1$ and $t{+}2$ terms.}
Because $X^\ast_t$ appears as the lag-1 term of the $t{+}1$ conditional and the
lag-2 term of the $t{+}2$ conditional, a proposal for $X^\ast_t$ changes three
factors of the joint prior. The updates add all three contributions to the
acceptance ratio (``Fix~2'', e.g.\ \code{update\_params\_ar2.R} L62--95); omitting
the lag-2 term would break detailed balance and bias the lag-2 autocorrelation
away from $\gamma_2$.
\end{enumerate}

Finally, the prior-predictive simulators are duals of the same maps:
\code{makeTauAR2}, \code{makeZAR2}, and \code{makeKnotsAR2}
(\code{auxfunctions\_ar2.R} L153--193) draw a unit-variance stationary AR(2) path
$X^\ast$ via \code{simulate\_ar2\_standard} and then apply the corresponding inverse
copula / inv.\ probit. The \code{*ARP} variants generalise the AR(2) to a general
AR($p$), solving the discrete Lyapunov equation for the stationary state covariance
while preserving $\gamma_0=1$; the logic is identical.

\paragraph{Conclusion.}
The code implements ``transform $\to$ standard normal $\to$ autocorrelated Gaussian
prior $\to$ back to the skew-$t$ marginal'' exactly as written, with the three
parameters $(\tau, z, w)$ mapped through Gamma, Half-Normal, and Uniform copula
marginals respectively. The only transform-irrelevant defect is the unused
\code{dhn} normalising constant.

\end{document}
